% CHECKED DEDUCTION: full self-reading, two independent whole-proof audits, and parent review passed.
% Existing checked p=2 and p=37 arguments remain separate.
\subsection{Several even directions and the range \texorpdfstring{$p<1000$}{p<1000}}
\label{sec:multieven-poles}

The preceding pole argument extends to several even arithmetic directions and to characteristic polynomials with repeated roots. The arithmetic hypothesis is most cleanly stated directly in Galois cohomology.

\begin{theorem}[K17-POLE-GENERAL]\label{thm:pole-general}
Let $p$ be an odd prime satisfying Vandiver, let $F=\mathbb Q(\mu_p)$, and let $P$ be the Galois group of its maximal pro-$p$ extension unramified outside $p$. Assume that the cup products of the odd-character summand of $H^1(P,\mathbb F_p)$ span $H^2(P,\mathbb F_p)$. Then
\[
 \Omega_p^*=\Omega_p.
\]
No assumption that the relevant Iwasawa polynomials have degree one, split over $\mathbb Q_p$, or have distinct roots is required.
\end{theorem}

Here odd means the $-1$ eigenspace for complex conjugation, equivalently the sum of the odd Teichm\"uller characters. Cup/relator duality says that the hypothesis supplies a nonzero odd--odd commutator coefficient in each nonzero even relation eigenspace. This is a statement about the given arithmetic group and its full cup product; it does not require a chosen basis of cyclotomic units.

\paragraph{The arithmetic module and its companion bases.}
Set $K=\mathbb Q(\mu_{p^\infty})$, $H=\operatorname{Gal}(\Omega_p/K)$, $\Gamma=\operatorname{Gal}(K/F)=\langle\gamma\rangle$, $\Lambda=\mathbb Z_p[[\Gamma]]$, and $M=H^{\mathrm{ab}}$. Choose $\gamma$ commuting with a lift of $\Delta=\operatorname{Gal}(F/\mathbb Q)$ and with $\chi(\gamma)=u=1+p$, as above.

Iwasawa's theorem and Ferrero--Washington make $H$ free pro-$p$, so $M$ is $\mathbb Z_p$-torsion-free. Weak Leopoldt gives $\operatorname{rank}_{\Lambda}M=(p-1)/2$. Kummer duality makes $M^+$ a torsion $\Lambda$-module; it is finitely generated over $\mathbb Z_p$ because it is pseudo-isomorphic to $(X_\infty^-)^\iota(1)$, whose $\mu$-invariant is zero. Finite kernels and cokernels preserve this finite generation. Thus $M^+$ is finite free over $\mathbb Z_p$ \cite[Theorem~2]{Iwasawa1981}\cite[Theorems~3.3.3--3.3.4, Proposition~3.4.7, Corollary~3.4.8]{SharifiIwasawaNotes}.

Under Vandiver the minimal generators of $P$ in McCallum--Sharifi, \S8, pp.~22--23, are $\gamma$, one $x_a$ for each odd residue class $a$, and one $y_e$ for each occurring even character $e$ \cite{McCallumSharifi2003}. From
\[
 0\longrightarrow M_\Gamma\longrightarrow P^{\mathrm{ab}}
  \longrightarrow\Gamma\longrightarrow0
\]
and Nakayama, $M^-$ has $(p-1)/2$ generators over $\Lambda$ and the same rank. The kernel of the resulting surjection from a free module has rank zero and is zero. Hence
\[
 M^-\simeq\bigoplus_{a\text{ odd}\bmod p-1}\Lambda e_a.       \tag{G.1}
\]
Each nonzero even component $T_e=M^{(e)}$ is cyclic over $\Lambda$ and finite free over $\mathbb Z_p$. Put $\lambda_e=\operatorname{rank}_{\mathbb Z_p}T_e$ and choose a cyclic generator $v_e$. If $q_e(z)$ is the characteristic polynomial of $\gamma$ on $T_e$, then
\[
 q_e(z)\equiv(z-1)^{\lambda_e}\pmod p.                         \tag{G.2}
\]
Indeed the continuous action of the pro-$p$ group $\Gamma$ on $T_e/pT_e$ has finite $p$-group image and is unipotent.

The vectors $v_e,\gamma v_e,\ldots,\gamma^{\lambda_e-1}v_e$ form a $\mathbb Z_p$-basis. Cayley--Hamilton makes their span stable under $\gamma$ and $\gamma^{-1}$, since $q_e$ has unit constant coefficient. The span is compact and closed, so contains the entire $\Lambda$-span of $v_e$. It therefore surjects onto $T_e$, and a surjection between free $\mathbb Z_p$-modules of the same finite rank is an isomorphism.

Choose actual lifts $x_a,y_e\in H$ and set
\[
 x_{a,t}=\gamma^t x_a\gamma^{-t},\quad t\in\mathbb Z_p,
 \qquad y_{e,j}=\gamma^j y_e\gamma^{-j},\quad0\le j<\lambda_e.
\]
These constitute a free compact basis of $H$. To recall the complete argument, take the inverse limit of the free pro-$p$ groups on the finite quotients of this compact generating space. Its abelianization is precisely (G.1) together with the indicated finite bases of the $T_e$. The natural map to $H$ is an isomorphism on abelianizations, hence is onto by the Frattini criterion for closed subgroups. Projectivity of $H$ supplies a section; its closed image is Frattini-surjective in the source and hence equals the source. The map is an isomorphism.

We may make the lifts exact nonlinear $\Delta$-eigenlifts by Sharifi's averaging lemma \cite[Lemma~2.1, pp.~3--4]{Sharifi2002}:
\[
 \delta x_a\delta^{-1}=x_a^{\omega(\delta)^a},\qquad
 \delta y_e\delta^{-1}=y_e^{\omega(\delta)^e}.
\]
This preserves the prescribed abelian eigenspace classes. Commutation of $\gamma$ and $\Delta$ gives the same eigenproperty for all their indicated conjugates.

\paragraph{The companion error and its cup coefficient.}
Write
\[
 q_e(z)=z^{\lambda_e}-\sum_{j=0}^{\lambda_e-1}c_{e,j}z^j
\]
and, with the product taken in increasing $j$, define
\[
 E_e=(\gamma^{\lambda_e}y_e\gamma^{-\lambda_e})
          \left(\prod_jy_{e,j}^{c_{e,j}}\right)^{-1}.
\]
The characteristic-polynomial relation in $M$ puts $E_e$ in the closed commutator subgroup of $H$. Its quadratic Magnus coefficients in two odd orbit letters give bounded measures and integral Amice transforms $C_{e;a,b}(T_1,T_2)$. We claim that for each $e$, there is an ordered odd pair $(a_e,b_e)$ such that
\[
 C_{e;a_e,b_e}(0,0)\in\mathbb Z_p^\times.                       \tag{G.3}
\]
The following relation-space argument supplies this comparison even when there are several even characters.

Let $F_{\mathrm{free}}$ be free pro-$p$ on $\gamma,x_a,y_e$, retaining only $j=0$ as an even generator. Substitute the actual conjugate words for all $y_{e,j}$ and $x_{a,t}$ in the compact free expansion of $E_e$. The group $P=H\rtimes\Gamma$ has one defining relator per even character:
\[
 R_e=\gamma^{\lambda_e}y_e\gamma^{-\lambda_e}
       \left(\prod_jy_{e,j}^{c_{e,j}}\right)^{-1}E_e^{-1}.       \tag{G.4}
\]
This is verified by the universal property. An assignment of $\gamma,x_a,y_e$ into a pro-$p$ target satisfying (G.4) gives images of all compact generators of $H$. Compatibility with $\gamma$ is automatic on the odd orbits and along each finite even chain except at its last element; the last compatibility is exactly (G.4). It therefore extends uniquely to the semidirect product.

Let $R$ be the closed relation subgroup. The exact eigenlifts let $\Delta$ act on $F_{\mathrm{free}}$ by fixing $\gamma$ and raising each other generator to its prescribed Teichm\"uller unit power. These are finite-order automorphisms; the quotient to $P$ is equivariant. Abelianize the kernel of $F_{\mathrm{free}}\to\Gamma$. Its free relative abelianization is
\[
 A_{\mathrm{free}}=
   \bigoplus_a\Lambda X_a\oplus\bigoplus_e\Lambda Y_e.
\]
The image of $R_e$ is exactly $q_e(\gamma)Y_e$. The word $E_e$ is a closed commutator word in the compact free basis, so has zero relative abelianization before imposing any arithmetic relations. Normal conjugation acts on this module through $\Gamma$. Hence the whole closed image is
\[
 R_{\mathrm{lin}}=\bigoplus_eq_e(\gamma)\Lambda Y_e,             \tag{G.5}
\]
a free $\Lambda$-module, with its indicated character $e$ on each summand.

Linearization gives a surjection
\[
 R/\overline{R^p[R,F_{\mathrm{free}}]}
       \longrightarrow R_{\mathrm{lin}}/(p,T)R_{\mathrm{lin}}.  \tag{G.6}
\]
The source is generated by the number of defining relators and the target has exactly that dimension. Thus (G.6) is an equivariant isomorphism, and the class of $R_e$ has character $e$. The cup-normalized character-$e$ relation is therefore a nonzero scalar multiple of $R_e$ in this quotient. All relators lie in the Frattini subgroup; passing to descending $p$-central degree two annihilates $R^p$ and $[R,F_{\mathrm{free}}]$.

The companion product in (G.4) contains only $y_e$ and its $\gamma$-conjugates and contributes no odd--odd bracket in that degree. Its odd--odd coefficient is consequently minus the augmentation coefficient of $E_e$, namely $-C_{e;a,b}(0,0)$ up to the commutator convention. The assumed odd cup surjectivity and the usual minimal-presentation cup/relator duality prove (G.3) \cite[\S8, pp.~21--23]{McCallumSharifi2003}. No freeness or asphericity claim about the full relation module is being used.

\paragraph{Quadratic target coefficients and partial fractions.}
Use the actual continuous dense rational representation $H\to U(\mathbb Q_p)$ and the completed free odd Lie algebra supplied by Hain--Matsumoto and Brown \cite{HainMatsumoto2001,Brown2019}. Conjugate rationally so that $\gamma$ is in a Levi torus; the successive positive-weight equations are solvable because $u^w-1\ne0$. The commuting $\Delta$ then also acts by its weight characters. Normalize each odd free letter $s_m$ so that the linear coefficient of $\rho(x_{a,t})-1$ on its matching residue class is $u^{mt}$. The necessary scalars are nonzero by density. Every $y_{e,j}$ has target generator length at least two by parity.

For an ordered pair $(m,n)$ put $z=u^{m+n}$, and let $B_{e,j;m,n}$ be the coefficient of $s_ms_n$ in $\rho(y_{e,j})-1$. The actual conjugate definition gives
\[
 B_{e,j;m,n}=z^jB_{e,0;m,n},
 \qquad
 q_e(z)B_{e,0;m,n}=C_{e;a,b}(u^m-1,u^n-1),                    \tag{G.7}
\]
where $a\equiv m$, $b\equiv n\pmod{p-1}$. The second equality is the length-two component of $E_e$. These coefficients have character $e$, so vanish when $a+b\not\equiv e$. For those residue classes (G.7) makes $C_{e;a,b}$ vanish on the entire Cartesian weight grid; the successive analytic identity theorem forces $C_{e;a,b}=0$ as a power series. Thus wrong-character numerators vanish identically, and the unit pair in (G.3) has $a_e+b_e\equiv e$.

Before dividing, exclude all exceptional weight sums. There are finitely many positive integers $w$ such that $u^w$ is a root of some $q_e$, since $u$ has infinite order. Restrict every test weight $m$ to exceed all such $w$. Then all pair denominators are nonzero. In each fixed odd residue class, the tails $m=a+(p-1)p^n$ still accumulate $p$-adically at $a$, and suffice for every analytic identity below.

Let $L/\mathbb Q_p$ be a finite extension splitting all $q_e$, with valuation ring $\mathcal O_L$ and uniformizer $\pi$. Every root $\alpha$ satisfies $\alpha-1\in\pi\mathcal O_L$ by (G.2). If its multiplicity is $n_{e,\alpha}$, ordinary partial fractions give the two bases of the same space of strictly proper rational functions:
\[
 \left\{\frac{z^j}{q_e(z)}:0\le j<\lambda_e\right\},
 \qquad
 \left\{\frac1{(z-\alpha)^k}:1\le k\le n_{e,\alpha}\right\}_{\alpha}.
\]
Apply the resulting constant invertible $L$-linear change to the degree-two Magnus letters $y_{e,j}$. This is a linear change in the associative coefficient algebra, not a change of actual pro-$p$ generators. At each fixed degree, only finitely many scalar coefficients occur, so measures retain a finite common denominator bound. The first nonzero assigned degree is preserved.

In the new letters $Y_{e,\alpha,k}$, the leading target kernels are
\[
 \frac{C_{e;a,b}(T_1,T_2)}
      {\bigl((1+T_1)(1+T_2)-\alpha\bigr)^k}.                   \tag{G.8}
\]
All even letters have assigned degree two and the orbit letters degree one. The measure-valued Magnus map remains injective after scalar extension and this invertible homogeneous change.

\paragraph{Colored pair tilings and Laurent coefficient separation.}
Suppose $g\ne1$ is in the kernel and let $d$ be its first nonzero assigned degree. A coefficient pattern is an ordered tiling of $d$ target positions by singleton odd letters and paired even letters. Each pair is colored by $(e,\alpha,k)$. For fixed target residue classes, its vanishing target coefficient gives
\[
 \sum_{\mathcal P} A_{\mathcal P}(T_i:i\text{ singleton})
 \prod_{(j,e,\alpha,k)\in\mathcal P}
 \frac{C_{e;a_j,a_{j+1}}(T_j,T_{j+1})}{D_{j,\alpha}^k}=0,
 \quad D_{j,\alpha}=(1+T_j)(1+T_{j+1})-\alpha.                \tag{G.9}
\]
All $A_{\mathcal P}$ are bounded $L$-valued measure transforms. Only the indicated first source degree contributes: a larger degree has larger target length, and every smaller degree is zero. Substitution at each finite arithmetic weight cutoff is continuous and multiplicative by integration of bounded continuous coefficient functions and uniform locally constant approximation, as in the preceding complete proofs.

There are finitely many patterns and denominators. Clearing denominators in (G.9) gives a bounded-coefficient analytic identity on the full Cartesian tails of the allowed weight grids. Successive one-variable identity theorems give an identity of power series over $L$. Its bounded-denominator ring is
\[
 R_L=\mathcal O_L[[T_1,\ldots,T_d]][1/\pi],
\]
not the ring permitting arbitrary unbounded coefficient denominators.

Choose a nonzero pattern with the maximum number of even letters. Fix its pair positions $J$ and its $(e,\alpha)$ colors. Among nonzero patterns with those positions and colors choose a pole-order vector $(k_j)$ maximal in the componentwise partial order. This is possible because there are finitely many patterns. Assign the unit target pair $(a_e,b_e)$ from (G.3) at each chosen pair and its source orbit type at each singleton.

For $j\in J$ introduce $S_j=D_{j,\alpha_j}$. The pairs are disjoint and the substitutions
\[
 T_{j+1}=\frac{\alpha_j+S_j}{1+T_j}-1
\]
give an isomorphism of formal $\mathcal O_L$ power-series rings: their constant terms are topologically nilpotent. Retain the unpaired variables and one left variable from each pair as base variables. Let $K_0$ be the fraction field of their bounded-denominator power-series ring. Equation (G.9) embeds into Laurent series over $K_0$ in the $S_j$, with finite lower bounds on all exponents.

Every other denominator is analytic in these variables. At the same pair position a different root gives $S_j+(\alpha_j-\alpha')$ with nonzero constant in $K_0$. At a different position its two entries lie in different blocks, whose remaining base variables are independent; their product is not identically its root. Its constant term is again nonzero in $K_0$. These denominators are therefore invertible in $K_0[[S_j:j\in J]]$. A different even character at a selected position contributes zero for the assigned target types, since their sum fixes $e$. Each selected $C$ remains a nonzero unit at $S=0$.

Extract the coefficient of $\prod_{j\in J}S_j^{-k_j}$. A pattern missing a selected pair or using a different root there has no negative exponent in that variable. A pattern properly containing $J$ has more even letters and zero coefficient by maximality. For patterns with the same positions and root colors, an order vector $k'$ can contribute only if $k'_j\ge k_j$ for every $j$, since all remaining series have nonnegative exponents. Componentwise maximality excludes every strict such vector. Thus only the selected pattern contributes, and its coefficient is
\[
 A_{\mathcal P}\prod_{j\in J}
 C_{e_j;a_{e_j},b_{e_j}}\big|_{S=0}\ne0.
\]
The unpaired-variable coefficient ring embeds in $K_0$, and all factors are units. This is a contradiction. When there are no even pairs, the same calculation reduces to odd-moment separation. Hence the actual rational action is injective on $H$. Cyclotomic abelianization detects any kernel outside $H$, and the rational action is induced by the pro-$p$ outer action. Theorem~\ref{thm:pole-general} follows.

\paragraph{K17-CUP-PARITY: from $p$-units to odd Galois generators.}
Let $A_F$ be the $p$-class group and
\[
 E=\mathcal O_F[1/p]^\times/(\mathcal O_F[1/p]^\times)^p.
\]
For a $p$-unit $v$, $v/c(v)$ is a global unit with all complex conjugates of absolute value one, hence a root of unity. An elementary proof is that the monic polynomials of its positive powers have bounded degree and bounded integer coefficients, so only finitely many roots occur and two powers coincide. Thus $E^-$ is exactly the one-dimensional space generated by $\mu_p$; it is nonzero because $F$ does not contain $\mu_{p^2}$.

The Kummer sequence and Vandiver give
\[
 0\longrightarrow E\longrightarrow H^1(P,\mu_p)
  \longrightarrow A_F[p]\longrightarrow0,
 \qquad E^+=H^1(P,\mu_p)^+.
\]
There is no Brauer contribution, since $F$ is totally imaginary with just one omitted finite prime, and
\[
 H^2(P,\mu_p^{\otimes2})=A_F\otimes\mu_p
\]
is a plus space. Mixed parity cups therefore vanish. The alternating cup on the one-dimensional $E^-$ vanishes because $p$ is odd. Surjectivity of the cup on all $p$-units consequently implies surjectivity from $E^+\otimes E^+$. Untwisting gives precisely the odd $H^1(P,\mathbb F_p)$ cup hypothesis of Theorem~\ref{thm:pole-general}. Under cup/relator duality, $x_a$ has character $\omega^a$, its Kummer dual has character $\omega^{1-a}$, and a relation of character $\omega^e$ corresponds to twisted target character $\omega^{2-e}$. Thus the indices satisfy $a+b\equiv e\pmod{p-1}$, as used above.

\begin{corollary}[K17-KERNEL-LT1000]\label{cor:kernel-lt1000}
For every odd prime $p<1000$, one has $\Omega_p^*=\Omega_p$.
\end{corollary}
\begin{proof}
Buhler--Harvey verify Vandiver for every prime below $163{,}577{,}356$, in particular for this entire range \cite[abstract and \S1, p.~1]{BuhlerHarvey2011}. Sharifi's Theorem~1.2 and Corollary~5.9 prove that the $p$-unit cup pairing is surjective for every odd $p<1000$ \cite[pp.~2 and~27]{Sharifi2007}. The parity lemma supplies the exact odd-generator hypothesis, so Theorem~\ref{thm:pole-general} applies.
\end{proof}

The cited computational theorems are used as published results: Sharifi's Theorem~5.7 is an explicit Magma modular-symbol computation, and the Vandiver range is the published Buhler--Harvey verification. These computations were not rerun here. The deduction imposes no separate numerical $\lambda$-invariant or simple-root condition. For an arbitrary remaining prime, the explicit sufficient conditions are Vandiver and the stated odd cup surjectivity. This argument does not settle their validity for all primes, and it does not prove integral odd generation of the image Lie algebra.
